\documentclass[11pt]{article}

\usepackage[T1]{fontenc}
\usepackage[utf8]{inputenc}
\usepackage{lmodern}
\usepackage{amsmath,amssymb,amsfonts,amsthm}
\usepackage[margin=1in]{geometry}
\usepackage{microtype}
\usepackage{xcolor}
\usepackage[colorlinks=true,linkcolor=teal,urlcolor=teal]{hyperref}

\newcommand{\R}{\mathbb R}
\newcommand{\C}{\mathbb C}
\newcommand{\Hc}{\mathsf H}
\newcommand{\F}{\mathcal F_{\mathrm{orb}}}
\newcommand{\T}{\mathcal T_1}
\newcommand{\K}{\mathcal K}
\newcommand{\A}{\mathfrak A}
\newcommand{\M}{\mathfrak M}
\newcommand{\Tr}{\operatorname{Tr}}
\newcommand{\dd}{\,\mathrm d}
\newcommand{\sh}{\operatorname{sh}}
\newcommand{\N}{\mathsf N}

\newtheorem*{lemmaC}{Lemma C.3 (rephrased)}

\title{A Short Guide to Lemma C.3\\
\large The residual subinstrument of a covariant hybrid map}
\author{}
\date{}

\begin{document}
\maketitle

\begin{abstract}
Lemma C.3 describes every exactly covariant hybrid map as a convolution of
the translated input with one residual quantum subinstrument.  Informally,
the deterministic part of the classical output is
$\sqrt\gamma$ times the classical input, while a residual variable records
all additional classical error and its correlated quantum operation.  This
note explains measure-valued instruments, the approximate-identity
construction of the residual object, the proof of the convolution formula,
and the specialization to the Gaussian input used in Proposition C.4.
\end{abstract}

\section{The structural question}

Let
\[
  \mathfrak T=L^1(\Hc;\T(\F)),
  \qquad
  \A=C_0(\Hc;\K(\F)),
  \qquad
  \M=\A^*.
\]
An element $R\in\mathfrak T$ is a classical--quantum input: at the
classical value $x\in\Hc$, it carries the trace-class operator $R(x)$.
A positive element of $\M$ is an operator-valued Radon measure on the output
classical space.

Suppose
\[
  \Gamma:\mathfrak T\longrightarrow\M
\]
is completely positive, trace nonincreasing, and covariant under both
classical translations and oscillator displacements.  Classical covariance
has gain $\sqrt\gamma$: translating the input by $x$ translates the output
by $\sqrt\gamma x$.  The question is whether the part of the output that
remains after removing this forced translation can be described by one
object independent of $x$.  Lemma C.3 answers yes.

\section{Background: quantum instruments and subinstruments}

A quantum instrument assigns to every Borel event $E\subset\Hc$ a completely
positive map $\mathcal I_E$ on trace-class operators.  For an input state
$X$, the number
\[
  \Tr\mathcal I_E(X)
\]
is the probability that the classical outcome lies in $E$, and the positive
operator $\mathcal I_E(X)$ is the corresponding unnormalized posterior
quantum state.  Countable additivity in $E$ is required in trace norm.

A \emph{subinstrument} permits total probability below one:
\[
  \Tr\mathcal I_{\Hc}(X)\leq\Tr X
  \qquad(X\geq0).
\]
This is appropriate here because the weak-* limiting map obtained in Lemma
C.2 may lose mass at infinity.  It is convenient to package all branches as
a measure-valued completely positive map
\[
  \mathcal I:\T(\F)\longrightarrow
  C_0(\Hc;\K(\F))^*.
\]
The value $\mathcal I(X)(E)$ is exactly $\mathcal I_E(X)$.

\section{The expected convolution formula}

Imagine first that the input is concentrated at one classical point $x$ and
has quantum part $X$.  Translation covariance says that its output should be
a translate by $\sqrt\gamma x$ of the output obtained at $x=0$.  If
$\mathcal I(X)$ denotes this origin output, then
\[
  \Gamma(\delta_x X)(E)
  =\mathcal I(X)(E-\sqrt\gamma x).
\]
An $L^1$ input is a continuous mixture of such point inputs, so one expects
\begin{equation}\label{eq:convolution-intuition}
  [\Gamma(R)](E)
  =\int_{\Hc}
  [\mathcal I(R(x))](E-\sqrt\gamma x)\dd x.
\end{equation}

There is one technical obstruction: $\delta_xX$ is not an element of
$L^1(\Hc;\T(\F))$.  Therefore $\Gamma(\delta_0X)$ is not directly defined.
The proof replaces the point mass by shrinking $L^1$ bumps and extracts a
weak-* limit.

\section{Statement of Lemma C.3}

\begin{lemmaC}
Let $\Gamma:\mathfrak T\to\M$ be completely positive, trace
nonincreasing, and exactly covariant under the hybrid translation--Weyl
group.  There exists a completely positive, trace-nonincreasing
measure-valued map
\[
  \mathcal I:\T(\F)\longrightarrow C_0(\Hc;\K(\F))^*
\]
whose Borel branches $\mathcal I_E$ form a subinstrument and such that
\begin{equation}\label{eq:convolution}
  [\Gamma(R)](E)
  =\int_{\Hc}
  [\mathcal I(R(x))](E-\sqrt\gamma x)\dd x.
\end{equation}
Each branch is displacement covariant with amplitude gain $\sqrt\gamma$.
If $\Gamma$ is covariant under independent oscillator phase rotations, then
every branch is phase covariant as well.

For the input
$R(x)=\varphi_{0,\Sigma_p}(x)\Phi_0$, if
$\Xi=\Gamma(R)$, then for bounded measurable $a$ and bounded oscillator
observable $M$,
\begin{equation}\label{eq:test}
  \int a(y)\Tr[M\,\Xi(\dd y)]
  =\int_{\Hc}
    \left(\int_{\Hc}a(y)
      \varphi_{0,\gamma\Sigma_p}(y-r)\dd y\right)
    \Tr[M\,\mathcal I(\Phi_0)(\dd r)].
\end{equation}
\end{lemmaC}

Equation \eqref{eq:convolution} is an equality of trace-class-valued Radon
measures.  Equation \eqref{eq:test} is its weak form against a classical test
function and a quantum observable.

\section{Proof, step by step}

\subsection{Step 1: approximate an input at the origin}

Choose nonnegative functions $u_j\in C_c(\Hc)$ such that
\[
  \int_{\Hc}u_j(x)\dd x=1,
  \qquad
  \operatorname{supp}u_j\downarrow\{0\}.
\]
For $X\in\T(\F)$, define
\begin{equation}\label{eq:Ij}
  \mathcal I_j(X)=\Gamma(u_jX),
  \qquad (u_jX)(x)=u_j(x)X.
\end{equation}
The field $u_jX$ is a normalized bump around the origin with quantum state
$X$.  Each $\mathcal I_j$ is completely positive.  For $X\geq0$,
\[
  \Tr\mathcal I_j(X)
  \leq\|u_jX\|_1
  =\Tr X,
\]
so it is trace nonincreasing.

The sequence need not converge in norm.  Pass to adjoints
\[
  \mathcal I_j^*:
  C_0(\Hc;\K(\F))\longrightarrow\mathcal B(\F).
\]
They are completely positive contractions.  Separability and a diagonal
weak-* extraction on a countable dense operator system give a pointwise
weak-* limit $\mathcal I^*$.  Because every matrix positive cone is weak-*
closed, $\mathcal I^*$ is completely positive.  Define $\mathcal I$ by
duality.

If $(e_\ell)$ is a contractive approximate identity of the test algebra,
then for $X\geq0$,
\[
  \Tr\mathcal I(X)
  =\sup_\ell\langle X,\mathcal I^*(e_\ell)\rangle
  \leq\Tr X.
\]
Thus the limit is still trace nonincreasing.  This is the residual
measure-valued map.

\subsection{Step 2: prove the convolution identity on simple tensors}

First take
\[
  R(x)=f(x)X,
  \qquad f\in L^1(\Hc),\quad X\in\T(\F).
\]
Let
\[
  (\sh_xu_j)(v)=u_j(v-x).
\]
Since
\[
  (f*u_j)(v)X
  =\int_{\Hc}f(x)(\sh_xu_j)(v)X\dd x,
\]
bounded linearity, Fubini's theorem, and classical translation covariance
give
\begin{equation}\label{eq:before-limit}
  \Gamma((f*u_j)X)
  =\int_{\Hc}f(x)\,
    \sh_{\sqrt\gamma x}\mathcal I_j(X)\dd x,
\end{equation}
where
\[
  (\sh_s\Xi)(E)=\Xi(E-s).
\]

The left side converges in measure norm because $f*u_j\to f$ in $L^1$ and
$\Gamma$ is bounded.  To pass to the limit on the right, pair with any
$A\in C_0(\Hc;\K(\F))$.  For each $x$, weak-* convergence of
$\mathcal I_j(X)$ gives convergence of the pairing.  The integrand is
bounded in absolute value by
\[
  |f(x)|\,\|X\|_1\,\|A\|,
\]
which is integrable.  Dominated convergence therefore changes
$\mathcal I_j$ to $\mathcal I$ in \eqref{eq:before-limit}.  Uniqueness of the
representing Radon measure gives
\[
  [\Gamma(fX)](E)
  =\int f(x)[\mathcal I(X)](E-\sqrt\gamma x)\dd x.
\]

Finite sums of simple tensors are dense in
$L^1(\Hc;\T(\F))$.  Both sides are bounded by $\|R\|_1$, so the formula
extends by continuity to every $R$, proving \eqref{eq:convolution}.

\subsection{Step 3: verify that the Borel branches form a subinstrument}

Complete positivity of the measure-valued map means that, whenever
$[X_{kl}]\geq0$, the matrix
\[
  [\mathcal I(X_{kl})]
\]
is a positive matrix-valued measure.  Approximate the indicator of a Borel
set $E$ by nonnegative continuous functions.  Evaluation on $E$ preserves
matrix positivity, hence
\[
  [\mathcal I_E(X_{kl})]\geq0.
\]
Therefore every branch $\mathcal I_E$ is completely positive.

For $X\geq0$, positivity of the measure implies
\[
  0\leq\mathcal I_E(X)\leq\mathcal I_{\Hc}(X),
\]
and consequently
\[
  \Tr\mathcal I_E(X)\leq\Tr X.
\]
Thus every branch is trace nonincreasing and bounded.  For disjoint events,
countable additivity holds in trace norm on positive inputs because the
increments are positive and their traces add.  Decomposing a general
trace-class operator into linear combinations of positive ones extends the
claim to all inputs.

\subsection{Step 4: pass quantum covariance to every branch}

Let $D(w)$ be a Weyl displacement.  The full covariance of $\Gamma$ implies,
already for the approximants,
\[
  \mathcal I_j(D(w)XD(w)^*)
  =D(\sqrt\gamma w)\mathcal I_j(X)D(\sqrt\gamma w)^*.
\]
Here the equality is one of operator-valued measures; the classical event is
not moved by a purely quantum displacement.  Passing to the weak-* limit and
evaluating on $E$ gives
\[
  \mathcal I_E(D(w)XD(w)^*)
  =D(\sqrt\gamma w)\mathcal I_E(X)D(\sqrt\gamma w)^*.
\]
The same argument applies to independent phase rotations.  This branchwise
covariance is essential: Proposition C.4 applies a quantum stochastic-order
bound separately to every residual event $E$.

\subsection{Step 5: specialize to the Gaussian input}

Set
\[
  R(x)=\varphi_{0,\Sigma_p}(x)\Phi_0.
\]
The convolution formula says that one may generate the output by the
following formal procedure:
\begin{enumerate}
  \item draw $X\sim\N(0,\Sigma_p)$;
  \item draw the residual outcome and quantum output according to
  $\mathcal I(\Phi_0)$;
  \item report the classical value $Y=\sqrt\gamma X+r$.
\end{enumerate}
Since $\sqrt\gamma X\sim\N(0,\gamma\Sigma_p)$, conditioning on $r$ gives
the density
\[
  y\longmapsto\varphi_{0,\gamma\Sigma_p}(y-r).
\]
Testing this measure against $a(y)M$ yields exactly \eqref{eq:test}.  The
calculation is initially justified for continuous compactly supported $a$;
a monotone-class argument extends it to bounded measurable $a$.

\section{A classical analogy}

Ignore the quantum component and suppose $K$ is a Markov kernel satisfying
\[
  K(E\mid x+s)=K(E-\sqrt\gamma s\mid x).
\]
Taking $s=-x$ gives
\[
  K(E\mid x)=K(E-\sqrt\gamma x\mid0).
\]
Thus every such kernel has the form
\[
  Y=\sqrt\gamma X+R,
\]
where the residual law of $R$ is independent of $X$.  Lemma C.3 is the
classical--quantum version: the residual outcome $r$ may be correlated with
a quantum operation, so its law and posterior state must be encoded by an
instrument rather than an ordinary probability measure.

\section{Dependencies and role in the paper}

The construction uses weak-* compactness in much the same way as Lemma C.2,
but for a localized family $\Gamma(u_jX)$ rather than for group-averaged
channels.  It also uses:
\begin{itemize}
  \item approximate identities in $L^1(\Hc)$;
  \item covariance and convolution;
  \item density of finite sums $f(x)X$ in the Bochner $L^1$ space;
  \item Radon-measure regularity and trace-norm countable additivity; and
  \item dominated convergence after testing against $\A$.
\end{itemize}

In Proposition C.4, the residual subinstrument isolates all deviations from
the ideal classical dilation.  Every branch has precisely the displacement
and phase covariance needed for Lemma D.4.  The Gaussian specialization
\eqref{eq:test} then separates the final estimate into a classical Gaussian
overlap and a quantum thermal-state overlap.  This is the source of the
product
\[
  B_{\mathrm{cl}}(\gamma,d)B_{\mathrm{orb}}(\gamma,p).
\]

\section{Minimal prerequisite checklist}

At an intuitive level, the lemma says only that translation covariance
forces an additive-noise representation.  A rigorous reading additionally
requires operator-valued Radon measures, completely positive instruments,
weak-* compactness of adjoints, approximate identities, Bochner $L^1$
spaces, and trace-norm countable additivity.  The approximate identity is not
an optional complication: it is what makes sense of ``the output at a
classical point'' when the original domain contains only $L^1$ densities.

\end{document}
